\documentclass[../../main.tex]{subfiles}% 注意这里的文件路径不能用 ./main.tex ,否则用latexmk编译子文件会报错
\graphicspath{{\subfix{./image/}}{\subfix{../../image/}}} % 指定图片引用路径,后续可以直接使用图片文件名
% 如果图片存放在上级目录,则使用 ../../image/ 作为备用路径

% 例如:
% \begin{figure}[H]
% \centering
% \includegraphics[width=0.3\textwidth]{图.png}
% \caption{}
% \label{figure:图}
% \end{figure}
% 注意:上述\label{}一定要放在\caption{}之后,否则引用图片序号会只会显示??.

\begin{document}

\section{Cauchy-Riemann方程}

\begin{definition}
设 \( f(z) = u(x,y) + \text{i}v(x,y) \) 是定义在区域 \( D \) 上的函数，\( z_0 = x_0 + \text{i}y_0 \in D \)。我们说 \( f \) 在 \( z_0 \) 处\textbf{实可微}，是指 \( u \) 和 \( v \) 作为 \( x,y \) 的二元函数在 \( (x_0,y_0) \) 处可微.
\end{definition}

\begin{theorem}
如果 $f(z)=u(x,y)+\mathrm{i}v(x,y)$ 在 $z_0=x_0+\mathrm{i}y_0$ 处实可微, 则有极限
\begin{align*}
\lim_{\Delta x\rightarrow 0} \frac{f(z_0+\Delta x)-f(z_0)}{\Delta x}=\frac{\partial u}{\partial x}(z_0)+\mathrm{i}\frac{\partial v}{\partial x}(z_0)=\frac{\partial u}{\partial x}(x_0,y_0)+\mathrm{i}\frac{\partial v}{\partial x}(x_0,y_0),
\end{align*}
\begin{align*}
\lim_{\Delta y\rightarrow 0} \frac{f(z_0+\mathrm{i}\Delta y)-f(z_0)}{\Delta y}=\frac{\partial u}{\partial y}(z_0)+\mathrm{i}\frac{\partial v}{\partial y}(z_0)=\frac{\partial u}{\partial y}(x_0,y_0)+\mathrm{i}\frac{\partial v}{\partial y}(x_0,y_0),
\end{align*}
分别记为 $\frac{\partial f}{\partial x}(z_0)=\frac{\partial f}{\partial x}(x_0,y_0),\frac{\partial f}{\partial y}(z_0)=\frac{\partial f}{\partial y}(x_0,y_0)$,也简记为
\begin{align*}
\frac{\partial f}{\partial x}=\frac{\partial u}{\partial x}+\mathrm{i}\frac{\partial v}{\partial x},\quad \frac{\partial f}{\partial y}=\frac{\partial u}{\partial y}+\mathrm{i}\frac{\partial v}{\partial y}.
\end{align*}
称为$f$的\textbf{一阶偏导数}.与实变函数中的偏导数一样，可同理定义$f$的$k$\textbf{阶偏导数}.
\end{theorem}
\begin{proof}

\end{proof}

\begin{definition}
设$z=x+\mathrm{i}y$是一个复数,把 \( z,\bar{z} \) 看成独立变量，定义算子
\begin{align}\label{eq:-----2.2.3}
\frac{\partial}{\partial z} = \frac{1}{2}\left( \frac{\partial}{\partial x} - \text{i}\frac{\partial}{\partial y} \right), \quad
\frac{\partial}{\partial \overline{z}} = \frac{1}{2}\left( \frac{\partial}{\partial x} + \text{i}\frac{\partial}{\partial y} \right).
\end{align}
\end{definition}
\begin{remark}
为什么要像\eqref{eq:-----2.2.3}式那样来定义算子 \( \frac{\partial}{\partial z} \) 和 \( \frac{\partial}{\partial \overline{z}} \) 呢？这是因为如果把复变函数 \( f(z) \) 写成
\begin{align*}
f(x,y) = f\left( \frac{z + \overline{z}}{2}, -\text{i}\frac{z - \overline{z}}{2} \right),
\end{align*}
把 \( z,\overline{z} \) 看成独立变量，分别对 \( z \) 和 \( \overline{z} \) 求偏导数，则得
\begin{align*}
\frac{\partial f}{\partial z} &= \frac{\partial f}{\partial x}\frac{\partial x}{\partial z} + \frac{\partial f}{\partial y}\frac{\partial y}{\partial z} = \frac{1}{2}\left( \frac{\partial f}{\partial x} - \text{i}\frac{\partial f}{\partial y} \right), \\
\frac{\partial f}{\partial \overline{z}} &= \frac{\partial f}{\partial x}\frac{\partial x}{\partial \overline{z}} + \frac{\partial f}{\partial y}\frac{\partial y}{\partial \overline{z}} = \frac{1}{2}\left( \frac{\partial f}{\partial x} + \text{i}\frac{\partial f}{\partial y} \right).
\end{align*}
这就是表达式\eqref{eq:-----2.2.3}的来源。这说明在进行微分运算时，可以把 \( z,\overline{z} \) 看成独立的变量。
\end{remark}

\begin{proposition}
设 \( f:D \to \mathbb{C} \) 是定义在区域 \( D \) 上的函数，\( z_0 \in D \)，那么 \( f \) 在 \( z_0 \) 处实可微的充分必要条件是
\begin{align}\label{equation::::1-4}
f(z_0 + \Delta z) - f(z_0) = \frac{\partial f}{\partial z}(z_0)\Delta z + \frac{\partial f}{\partial \overline{z}}(z_0)\overline{\Delta z} + o(|\Delta z|).
\end{align}
成立.
\end{proposition}
\begin{proof}
设 \( f \) 在 \( z_0 \) 处实可微，由二元实值函数可微的定义，有
\begin{align}
u(x_0 + \Delta x, y_0 + \Delta y) - u(x_0, y_0) &= \frac{\partial u}{\partial x}(x_0, y_0)\Delta x + \frac{\partial u}{\partial y}(x_0, y_0)\Delta y + o(|\Delta z|), \label{eq:-----2.2.1} \\
v(x_0 + \Delta x, y_0 + \Delta y) - v(x_0, y_0) &= \frac{\partial v}{\partial x}(x_0, y_0)\Delta x + \frac{\partial v}{\partial y}(x_0, y_0)\Delta y + o(|\Delta z|), \label{eq:-----2.2.2}
\end{align}
这里，\( |\Delta z| = \sqrt{(\Delta x)^2 + (\Delta y)^2} \)。于是
\begin{align*}
f(z_0 + \Delta z) - f(z_0) &= u(x_0 + \Delta x, y_0 + \Delta y) - u(x_0, y_0) + \text{i}(v(x_0 + \Delta x, y_0 + \Delta y) - v(x_0, y_0)) \\
&= \frac{\partial u}{\partial x}(x_0, y_0)\Delta x + \frac{\partial u}{\partial y}(x_0, y_0)\Delta y + o(|\Delta z|) + \text{i}\left( \frac{\partial v}{\partial x}(x_0, y_0)\Delta x + \frac{\partial v}{\partial y}(x_0, y_0)\Delta y + o(|\Delta z|) \right) \\
&= \left( \frac{\partial u}{\partial x}(x_0, y_0) + \text{i}\frac{\partial v}{\partial x}(x_0, y_0) \right)\Delta x + \left( \frac{\partial u}{\partial y}(x_0, y_0) + \text{i}\frac{\partial v}{\partial y}(x_0, y_0) \right)\Delta y + o(|\Delta z|) \\
&= \frac{\partial f}{\partial x}(x_0, y_0)\Delta x + \frac{\partial f}{\partial y}(x_0, y_0)\Delta y + o(|\Delta z|).
\end{align*}
把 \( \Delta x = \frac{1}{2}(\Delta z + \overline{\Delta z}) \)，\( \Delta y = \frac{1}{2\text{i}}(\Delta z - \overline{\Delta z}) \) 代入上式，得
\begin{align*}
f(z_0 + \Delta z) - f(z_0) &= \frac{1}{2}\frac{\partial f}{\partial x}(x_0, y_0)(\Delta z + \overline{\Delta z}) - \frac{\text{i}}{2}\frac{\partial f}{\partial y}(x_0, y_0)(\Delta z - \overline{\Delta z}) + o(|\Delta z|) \\
&= \frac{1}{2}\left( \frac{\partial}{\partial x} - \text{i}\frac{\partial}{\partial y} \right)f(x_0, y_0)\Delta z + \frac{1}{2}\left( \frac{\partial}{\partial x} + \text{i}\frac{\partial}{\partial y} \right)f(x_0, y_0)\overline{\Delta z} + o(|\Delta z|).
\end{align*}
从而上式可写为
\begin{align}\label{eq:-----2.2.4}
f(z_0 + \Delta z) - f(z_0) &= \frac{\partial f}{\partial z}(z_0)\Delta z + \frac{\partial f}{\partial \overline{z}}(z_0)\overline{\Delta z} + o(|\Delta z|).
\end{align}
容易看出\eqref{eq:-----2.2.4}式和\eqref{eq:-----2.2.1}\eqref{eq:-----2.2.2}两式等价。
\end{proof}

\begin{theorem}\label{theorem:复变函数可微的充要条件1}
设 \( f \) 是定义在区域 \( D \) 上的函数，\( z_0 \in D \)，那么 \( f \) 在 \( z_0 \) 处可微的充要条件是 \( f \) 在 \( z_0 \) 处实可微且 \( \frac{\partial f}{\partial \overline{z}}(z_0) = 0 \)。并且$f$在可微的情况下，\( f'(z_0) = \frac{\partial f}{\partial z}(z_0) \)。
\end{theorem}
\begin{proof}
如果 \( f \) 在 \( z_0 \) 处可微，由\eqref{eq:2.1.2}式得
\[
f(z_0 + \Delta z) - f(z_0) = f'(z_0)\Delta z + o(|\Delta z|)
\]
与\eqref{equation::::1-4}式比较就知道，\( f \) 在 \( z_0 \) 处是实可微的，而且 \( \frac{\partial f}{\partial \overline{z}}(z_0) = 0 \)，\( f'(z_0) = \frac{\partial f}{\partial z}(z_0) \)。

反之，若 \( f \) 在 \( z_0 \) 处实可微，且 \( \frac{\partial f}{\partial \overline{z}}(z_0) = 0 \)，则由\eqref{equation::::1-4}式得
\[
f(z_0 + \Delta z) - f(z_0) = \frac{\partial f}{\partial z}(z_0)\Delta z + o(|\Delta z|)
\]
由此即知 
\begin{align*}
\underset{\Delta z}{\lim}\frac{f\left( z_0+\Delta z \right) -f\left( z_0 \right)}{\Delta z}=\frac{\partial f}{\partial z}\left( z_0 \right).
\end{align*}
故\( f \) 在 \( z_0 \) 处可微，而且 \( f'(z_0) = \frac{\partial f}{\partial z}(z_0) \)。
\end{proof}

\begin{definition}[Cauchy-Riemann 方程]
设$f$是定义在区域$D$上的函数,$\frac{\partial f}{\partial \overline{z}} = 0$ 称为 $\mathbf{Cauchy}-\mathbf{Riemann}$\textbf{方程}.简称$\mathbf{C}.-\mathbf{R}.$\textbf{方程}
\end{definition}
\begin{note}
从这个方程可以得到 $f$ 的实部和虚部应满足的条件。
\end{note}

\begin{theorem}[Cauchy-Riemann方程的等价定义]\label{theorem:Cauchy-Riemann方程的等价定义}
设 $z=x+\mathrm{i}y$,$f(z) = u(x,y) + \mathrm{i}v(x,y)$,则

(i)Cauchy-Riemann方程$\frac{\partial f}{\partial \overline{z}} = 0$ 等价于
\begin{align}\label{equation:--2.5}
\begin{cases}
\frac{\partial u}{\partial x}=\frac{\partial v}{\partial y},\\
\frac{\partial u}{\partial y}=-\frac{\partial v}{\partial x}.
\end{cases}
\end{align}

(ii)Cauchy-Riemann方程$\frac{\partial f}{\partial \overline{z}} = 0$ 等价于
$$
\frac{\partial \overline{f}}{\partial z}=0.
$$

(iii)令$x=r\cos \theta ,y=r\sin \theta $,进而$z=r(\cos\theta+\mathrm{i}\sin\theta),f(z) = u(r,\theta) + \mathrm{i}v(r,\theta)$,则Cauchy-Riemann方程$\frac{\partial f}{\partial \overline{z}} = 0$ 等价于
\begin{align*}
\begin{cases}
\dfrac{\partial u}{\partial r} = \dfrac{1}{r} \dfrac{\partial v}{\partial \theta}, \\
\dfrac{\partial v}{\partial r} = -\dfrac{1}{r} \dfrac{\partial u}{\partial \theta}.
\end{cases}
\end{align*}
\end{theorem}
\begin{proof}
(i)由\eqref{eq:-----2.2.3}式得
\begin{align*}
\frac{\partial f}{\partial \overline{z}}=\frac{\partial u}{\partial \overline{z}} + \mathrm{i}\frac{\partial v}{\partial \overline{z}}=\frac{1}{2}\left(\frac{\partial u}{\partial x}+\mathrm{i}\frac{\partial u}{\partial y}\right)+\frac{\mathrm{i}}{2}\left(\frac{\partial v}{\partial x}+\mathrm{i}\frac{\partial v}{\partial y}\right)=\frac{1}{2}\left(\frac{\partial u}{\partial x}-\frac{\partial v}{\partial y}\right)+\frac{\mathrm{i}}{2}\left(\frac{\partial u}{\partial y}+\frac{\partial v}{\partial x}\right).
\end{align*}
因此Cauchy-Riemann 方程 $\frac{\partial f}{\partial \overline{z}} = 0$ 就等价于
\begin{align*}
\begin{cases}
\frac{\partial u}{\partial x}=\frac{\partial v}{\partial y},\\
\frac{\partial u}{\partial y}=-\frac{\partial v}{\partial x}.
\end{cases}
\end{align*}

(ii)又注意到
\begin{align*}
\frac{\partial \overline{f}}{\partial z}=\frac{\partial u}{\partial z}-\mathrm{i}\frac{\partial v}{\partial z}=\frac{1}{2}\left( \frac{\partial u}{\partial x}-\mathrm{i}\frac{\partial u}{\partial y} \right) -\frac{\mathrm{i}}{2}\left( \frac{\partial v}{\partial x}-\mathrm{i}\frac{\partial v}{\partial y} \right) =\frac{1}{2}\left( \frac{\partial u}{\partial x}-\frac{\partial v}{\partial y} \right) -\frac{\mathrm{i}}{2}\left( \frac{\partial u}{\partial y}+\frac{\partial v}{\partial x} \right) ,
\end{align*}
故Cauchy-Riemann 方程 $\frac{\partial f}{\partial \overline{z}} = 0$ 也等价于$\frac{\partial \overline{f}}{\partial z}=0.$

(iii)直接计算得
\begin{align*}
\frac{\partial \left( x,y \right)}{\partial \left( r,\theta \right)}=\begin{vmatrix}
\cos \theta& -r\sin \theta\\
\sin \theta& r\cos \theta\\
\end{vmatrix}=r.
\end{align*}
再由\refthe{Basis of Analytics-theorem:反函数组可微定理}知
\begin{align*}
\frac{\partial r}{\partial x}=\frac{\partial y}{\partial \theta}\bigg/\frac{\partial \left( x,y \right)}{\partial \left( r,\theta \right)}=\cos \theta,\quad \frac{\partial r}{\partial y}=-\frac{\partial x}{\partial \theta}\bigg/\frac{\partial \left( x,y \right)}{\partial \left( r,\theta \right)}=\sin \theta,\\
\frac{\partial \theta}{\partial x}=-\frac{\partial y}{\partial r}\bigg/\frac{\partial \left( x,y \right)}{\partial \left( r,\theta \right)}=-\frac{\sin \theta}{r},\quad \frac{\partial \theta}{\partial y}=\frac{\partial x}{\partial r}\bigg/\frac{\partial \left( x,y \right)}{\partial \left( r,\theta \right)}=\frac{\cos \theta}{r}.
\end{align*}
于是
\begin{gather*}
\frac{\partial u}{\partial x}=\frac{\partial u}{\partial r}\cdot \frac{\partial r}{\partial x}+\frac{\partial u}{\partial \theta}\cdot \frac{\partial \theta}{\partial x}=\frac{\partial u}{\partial r}\cdot \cos \theta -\frac{\partial u}{\partial \theta}\cdot \frac{\sin \theta}{r},
\\
\frac{\partial u}{\partial y}=\frac{\partial u}{\partial r}\cdot \frac{\partial r}{\partial y}+\frac{\partial u}{\partial \theta}\cdot \frac{\partial \theta}{\partial y}=\frac{\partial u}{\partial r}\cdot \sin \theta +\frac{\partial u}{\partial \theta}\cdot \frac{\cos \theta}{r},
\\
\frac{\partial v}{\partial x}=\frac{\partial v}{\partial r}\cdot \frac{\partial r}{\partial x}+\frac{\partial v}{\partial \theta}\cdot \frac{\partial \theta}{\partial x}=\frac{\partial v}{\partial r}\cdot \cos \theta -\frac{\partial v}{\partial \theta}\cdot \frac{\sin \theta}{r},
\\
\frac{\partial v}{\partial y}=\frac{\partial v}{\partial r}\cdot \frac{\partial r}{\partial y}+\frac{\partial v}{\partial \theta}\cdot \frac{\partial \theta}{\partial y}=\frac{\partial v}{\partial r}\cdot \sin \theta +\frac{\partial v}{\partial \theta}\cdot \frac{\cos \theta}{r}.
\end{gather*}
从而
\[
\frac{\partial u}{\partial x} = \frac{\partial v}{\partial y} \Longleftrightarrow \frac{\partial u}{\partial r} \cdot \cos\theta - \frac{\partial u}{\partial \theta} \cdot \frac{\sin\theta}{r} = \frac{\partial v}{\partial r} \cdot \sin\theta + \frac{\partial v}{\partial \theta} \cdot \frac{\cos\theta}{r}
\]
\[
\Longleftrightarrow \frac{\partial u}{\partial r} \cdot r\cos\theta - \frac{\partial u}{\partial \theta} \cdot \sin\theta = \frac{\partial v}{\partial r} \cdot r\sin\theta + \frac{\partial v}{\partial \theta} \cdot \cos\theta,
\]
\[
\frac{\partial u}{\partial y} = -\frac{\partial v}{\partial x} \Longleftrightarrow \frac{\partial u}{\partial r} \cdot \sin\theta + \frac{\partial u}{\partial \theta} \cdot \frac{\cos\theta}{r} = -\frac{\partial v}{\partial r} \cdot \cos\theta + \frac{\partial v}{\partial \theta} \cdot \frac{\sin\theta}{r}
\]
\[
\Longleftrightarrow \frac{\partial u}{\partial r} \cdot r\sin\theta + \frac{\partial u}{\partial \theta} \cdot \cos\theta = -\frac{\partial v}{\partial r} \cdot r\cos\theta + \frac{\partial v}{\partial \theta} \cdot \sin\theta.
\]
由 (i) 可知 Cauchy-Riemann 方程等价于 \( \frac{\partial u}{\partial x} = \frac{\partial v}{\partial y}, \frac{\partial u}{\partial y} = -\frac{\partial v}{\partial x} \)，故此时 Cauchy-Riemann 方程等价于
\[
\begin{cases}
\frac{\partial u}{\partial r} \cdot r\cos\theta - \frac{\partial u}{\partial \theta} \cdot \sin\theta = \frac{\partial v}{\partial r} \cdot r\sin\theta + \frac{\partial v}{\partial \theta} \cdot \cos\theta, \\
\frac{\partial u}{\partial r} \cdot r\sin\theta + \frac{\partial u}{\partial \theta} \cdot \cos\theta = -\frac{\partial v}{\partial r} \cdot r\cos\theta + \frac{\partial v}{\partial \theta} \cdot \sin\theta. \\
\end{cases}
\]
化简可得
\[
\begin{cases}
\frac{\partial u}{\partial r} = \frac{1}{r}\frac{\partial v}{\partial \theta}, \\
\frac{\partial v}{\partial r} = -\frac{1}{r}\frac{\partial u}{\partial \theta}. \\
\end{cases}
\]
\end{proof}

\begin{theorem}\label{theorem:复变函数可微的充要条件2}
设 \( f = u + \mathrm{i}v \) 是定义在区域 \( D \) 上的函数，\( z_0 = x_0 + \mathrm{i}y_0 \in D \)，那么 \( f \) 在 \( z_0 \) 处可微的充要条件是 \( u(x, y) \)，\( v(x, y) \) 在 \( (x_0, y_0) \) 处可微，且在 \( (x_0, y_0) \) 处满足Cauchy-Riemann方程,即
\[
\frac{\partial f}{\partial \overline{z}}=0\Longleftrightarrow \frac{\partial \overline{f}}{\partial z}=0\Longleftrightarrow \begin{cases}
\frac{\partial u}{\partial x}=\frac{\partial v}{\partial y},\\
\frac{\partial u}{\partial y}=-\frac{\partial v}{\partial x}\\
\end{cases}.
\]
在可微的情况下，有
\begin{align*}
f'(z_0) = \dfrac{\partial u}{\partial x} + \mathrm{i}\dfrac{\partial v}{\partial x}
= \dfrac{\partial v}{\partial y} + \mathrm{i}\dfrac{\partial v}{\partial x} = \dfrac{\partial u}{\partial x} - \mathrm{i}\dfrac{\partial u}{\partial y} = \dfrac{\partial v}{\partial y} - \mathrm{i}\dfrac{\partial u}{\partial y}.
\end{align*}
这里的偏导数都在 \( (x_0, y_0) \) 处取值。
\end{theorem}
\begin{proof}
由 \refthe{theorem:复变函数可微的充要条件1} 和 \hyperref[theorem:Cauchy-Riemann方程的等价定义]{Cauchy-Riemann方程的等价定义}容易得到充要条件的证明.再由 \refthe{theorem:复变函数可微的充要条件1}中的\( f'(z_0) = \dfrac{\partial f}{\partial z}(z_0) \) 和 \hyperref[theorem:Cauchy-Riemann方程的等价定义]{Cauchy-Riemann方程的等价定义}可得,在可微的情况下,有
\begin{align*}
f' (z_0)&=\frac{\partial f}{\partial z}(z_0)=\frac{1}{2}\left[ \frac{\partial f}{\partial x}(z_0)-\mathrm{i}\frac{\partial f}{\partial y}(z_0) \right] 
\\
&=\frac{1}{2}\left[ \frac{\partial u}{\partial x}(z_0)+\mathrm{i}\frac{\partial v}{\partial x}(z_0)-\mathrm{i}\left( \frac{\partial u}{\partial y}(z_0)+\mathrm{i}\frac{\partial v}{\partial y}(z_0) \right) \right] 
\\
&=\frac{1}{2}\left[ \frac{\partial u}{\partial x}(z_0)+\frac{\partial v}{\partial y}(z_0)+\mathrm{i}\left( \frac{\partial v}{\partial x}(z_0)-\frac{\partial u}{\partial y}(z_0) \right) \right] 
\\
&=\frac{1}{2}\left[ 2\frac{\partial u}{\partial x}(z_0)+2\mathrm{i}\frac{\partial v}{\partial x}(z_0) \right] 
\\
&=\frac{\partial u}{\partial x}(z_0)+\mathrm{i}\frac{\partial v}{\partial x}(z_0).
\end{align*}
再由\hyperref[theorem:Cauchy-Riemann方程的等价定义]{Cauchy-Riemann方程的等价定义}可得结论.
\end{proof}

\begin{definition}
\begin{enumerate}
\item 设 \( D \) 是 \( \mathbb{C} \) 中的区域，我们用 \( C(D) \) 记 \( D \) 上连续函数的全体，用 \( H(D) \) 记 \( D \) 上全纯函数的全体.

\item 设 \( f = u + \mathrm{i}v \)，我们用 \( C^1(D) \) 记有1阶连续偏导数的函数的全体。
用 \( C^k(D) \) 记在 \( D \) 上有 \( k \) 阶连续偏导数的函数的全体，\( C^\infty(D) \) 记在 \( D \) 上有任意阶连续偏导数的函数的全体。
\end{enumerate}
\end{definition}

\begin{proposition}\label{proposition:全纯函数必任意连续可微}
\begin{enumerate}[(1)]
\item \( H(D) \subseteq C(D) \).

\item $C^1(D) \subseteq C(D).$

\item 区域 \( D \) 上的全纯函数在 \( D \) 上有任意阶的连续偏导数，并且有如下的包含关系：
\[
H(D) \subseteq C^\infty(D) \subseteq C^k(D) \subseteq C^1(D) \subseteq C(D).
\]
这里，\( k \) 是大于 1 的自然数。
\end{enumerate}
\end{proposition}
\begin{proof}
\begin{enumerate}[(1)]
\item \refpro{proposition:全纯函数必连续}告诉我们，\( H(D) \subseteq C(D) \)。

\item 设 \( f = u + \mathrm{i}v \)，记 \( \frac{\partial f}{\partial x} = \frac{\partial u}{\partial x} + \mathrm{i}\frac{\partial v}{\partial x} \)，\( \frac{\partial f}{\partial y} = \frac{\partial u}{\partial y} + \mathrm{i}\frac{\partial v}{\partial y} \)。我们用 \( C^1(D) \) 记 \( \frac{\partial f}{\partial x} \)，\( \frac{\partial f}{\partial y} \) 在 \( D \) 上连续的 \( f \) 的全体。进而$u,v$关于$x,y$的偏导在$D$上都连续,由多元微积分的知识知道，$u,v$在$D$上都可微。于是对于任意 \( f \in C^1(D) \)，\( f \) 在 \( D \) 上实可微，从\eqref{eq:-----2.2.4}式知道
\begin{align*}
f(z_0 + \Delta z) - f(z_0) &= \frac{\partial f}{\partial z}(z_0)\Delta z + \frac{\partial f}{\partial \overline{z}}(z_0)\overline{\Delta z} + o(|\Delta z|).
\end{align*}
令$\Delta z\to 0$,则$\lim_{\Delta z\to 0}f(z_0+\Delta z)=f(z_0)$,
故\( f \) 在 \( D \) 上连续，因而
$C^1(D) \subseteq C(D).$

\item 
\end{enumerate}
\end{proof}

\begin{example}
研究函数 $f(z) = z^n$，$n$ 是自然数.
\end{example}
\begin{solution}
显然，$\frac{\partial f}{\partial \overline{z}} = 0$，且 $f$ 在整个平面上是实可微的. 因而，$f$ 是 $\mathbb{C}$ 上的全纯函数，而且
\begin{align}
f'(z) &= \frac{\partial f}{\partial z} = nz^{n - 1}. \label{eq:2.2.5}
\end{align}

\end{solution}

\begin{example}
研究函数 $f(z) = \mathrm{e}^{-\vert z \vert^2}$.
\end{example}
\begin{solution}
把 $f$ 写为 $f(z) = \mathrm{e}^{-z\overline{z}}$，于是 $\frac{\partial f}{\partial \overline{z}} = -\mathrm{e}^{-z\overline{z}}z$，它只有在 $z = 0$ 处才等于零. 因此，$\mathrm{e}^{-\vert z \vert^2}$ 只有在 $z = 0$ 处可微，它在任何点处都不是全纯的. 但它对 $x, y$ 有任意阶连续偏导数，所以它是 $C^\infty(\mathbb{C})$ 中的函数.

\end{solution}

\begin{proposition}\label{proposition:复变函数导数为0必是常函数}
设函数$f(z)=u+\mathrm{i}v$在区域$D$内解析,并满足下列条件之一:
\begin{enumerate}[(1)]
\item
\label{proposition:复变函数导数为0必是常函数-1} 在$D$内$f'(z)=0$.

\item\label{proposition:复变函数导数为0必是常函数-2} $\overline{f(z)}$在$D$内解析.

\item\label{proposition:复变函数导数为0必是常函数-3} $|f(z)|$在$D$内为常数.

\item\label{proposition:复变函数导数为0必是常函数-4} $\mathrm{Re}\,f(z)$或$\mathrm{Im}\,f(z)$在$D$内为常数.

\item\label{proposition:复变函数导数为0必是常函数-6} $\mathrm{Arg}f\left( z \right) $在$D$内为常数.

\item\label{proposition:复变函数导数为0必是常函数-5} $au+bv=c$,其中$a,b$与$c$是不全为零的实常数.
\end{enumerate}
则$f(z)$是常数.
\end{proposition}
\begin{proof}
设$f(z)=u(x,y)+\mathrm{i}v(x,y)$.
\begin{enumerate}[(1)]
\item 因为 \( f'(z) = 0 \)，所以由\refthe{theorem:复变函数可微的充要条件2}可知\( f'(z) = \frac{\partial u}{\partial x} + \mathrm{i}\frac{\partial v}{\partial x} = 0 \),并且\( \frac{\partial u}{\partial x} = \frac{\partial v}{\partial y} \)，\( \frac{\partial u}{\partial y} = -\frac{\partial v}{\partial x} \).于是$\frac{\partial u}{\partial x}=\frac{\partial u}{\partial y}=0,\frac{\partial u}{\partial y}=\frac{\partial v}{\partial y}=0$.
因此 \( u,v \)都是常数,故$f$是一常数.

\item 若$\overline{f(z)}=u(x,y)-\mathrm{i}v(x,y)$解析,由$\mathrm{C.-R.}$方程很容易得到$u_x=u_y=v_x=v_y=0$,从而$u(x,y),v(x,y)$在$D$内为常数,故$f(z)$亦为常数.

\item 若$|f(z)|\equiv C=0$,则显然$f(z)\equiv0$.
若$|f(z)|\equiv C\neq0$,则$f(z)\neq0$.

{\color{blue}证法一:}由$\overline{f(z)}\cdot f(z)\equiv C^2$,可推得$\overline{f(z)}\equiv\frac{C^2}{f(z)}$也是解析函数,
则利用\rrefpro{proposition:复变函数导数为0必是常函数}{proposition:复变函数导数为0必是常函数-2}可知$f(z)$在$D$内为常数.

{\color{blue}证法二:}由$|f(z)|^2=C^2$可得$u^2+v^2=C^2\neq0$,分别对$x,y$微分,再应用$\mathrm{C.-R.}$方程,讨论解二元一次齐次方程组,
即得在$D$内$u_x=v_y=u_y=v_x=0$.

\item 若$\mathrm{Re}\,f(z)$在$D$内为常数,可设$u(x,y)\equiv C$,则$u_x=u_y=0$,由$\mathrm{C.-R.}$方程,得
\begin{align*}
v_x=-u_y\equiv0, \quad u_x=v_y\equiv0.
\end{align*}
因此$u(x,y)\equiv C_1$,$v(x,y)\equiv C_2$,$f(z)=C_1+\mathrm{i}C_2$为常数.
若$\mathrm{Im}\,f(z)$在$D$内为常数,同上可得证.

\item 利用C.-R.方程易得.

\item 由条件可得$au_x+bv_x=0$,$\,\,au_y+bv_y=0$.

情况1:若$a\neq0$,$b=0$或$a=0$,$b\neq0$,则再由\rrefpro{proposition:复变函数导数为0必是常函数}{proposition:复变函数导数为0必是常函数-4}可知$u$或$v$为常数.

情况2:若$ab\neq0$,则根据
\begin{align*}
au_x-bu_y=0, \quad bu_x+au_y=0,
\end{align*}
注意到$\begin{vmatrix} a & -b \\ b & a \end{vmatrix}=a^2+b^2>0$,
系数矩阵满秩.
则上述方程组只有零解,故$u$为常数.同理$v$为常数.故$f(z)$是常数.
\end{enumerate}
\end{proof}

\begin{definition}[调和函数]
设 $u$ 是区域 $D$ 上的实值函数，如果 $u \in C^2(D)$，且对任意 $z \in D$，有
\begin{align}
\Delta u(z) &= \frac{\partial^2 u(z)}{\partial x^2} + \frac{\partial^2 u(z)}{\partial y^2} = 0, \label{eq:2.2.7}
\end{align}
就称 $u$ 是 $D$ 中的\textbf{调和函数}. $\Delta = \frac{\partial^2}{\partial x^2} + \frac{\partial^2}{\partial y^2}$ 称为 $\mathbf{Laplace}$\textbf{算子}.
\end{definition}

\begin{proposition}
设 $u \in C^2(D)$，那么 $\Delta u = 4 \frac{\partial^2 u}{\partial z \partial \overline{z}}$.
\end{proposition}
\begin{proof}
由\eqref{eq:-----2.2.3}式，有
\begin{align*}
\frac{\partial}{\partial z} &= \frac{1}{2}\left( \frac{\partial}{\partial x} - \text{i}\frac{\partial}{\partial y} \right), \\
\frac{\partial}{\partial \overline{z}} &= \frac{1}{2}\left( \frac{\partial}{\partial x} + \text{i}\frac{\partial}{\partial y} \right),
\end{align*}
所以
\begin{align*}
\frac{\partial^2 u}{\partial z \partial \overline{z}} &= \frac{\partial}{\partial \overline{z}} \frac{\partial u}{\partial z} = \frac{1}{4} \left[ \frac{\partial}{\partial x} \left( \frac{\partial u}{\partial x} - \mathrm{i} \frac{\partial u}{\partial y} \right) + \mathrm{i} \frac{\partial}{\partial y} \left( \frac{\partial u}{\partial x} - \mathrm{i} \frac{\partial u}{\partial y} \right) \right] \\
&= \frac{1}{4} \left( \frac{\partial^2 u}{\partial x^2} + \frac{\partial^2 u}{\partial y^2} \right) = \frac{1}{4} \Delta u.
\end{align*}
\end{proof}

\begin{theorem}\label{theorem:全纯函数的实部和虚部都是调和函数}
设 $f = u + \mathrm{i}v \in H(D)$，那么 $u$ 和 $v$ 都是 $D$ 上的调和函数.
\end{theorem}
\begin{proof}
因为 $f \in H(D)$，由\refthe{theorem:复变函数可微的充要条件2}，有
\begin{align*}
\frac{\partial f}{\partial \overline{z}} = 0, \quad
\frac{\partial \overline{f}}{\partial z} = 0.
\end{align*}
所以
\begin{align*}
\frac{\partial^2 f}{\partial z \partial \overline{z}} &= \frac{\partial^2 \overline{f}}{\partial z \partial \overline{z}} = 0.
\end{align*}
于是，由 $u = \frac{1}{2}(f + \overline{f})$ 即得
\begin{align*}
\Delta u &= 4 \frac{\partial^2 u}{\partial z \partial \overline{z}} = 0.
\end{align*}

同理可证 $\Delta v = 0$.
\end{proof}

\begin{definition}[共轭调和函数]
设 $u$ 和 $v$ 是 $D$ 上的一对调和函数，如果它们还满足 Cauchy-Riemann 方程
\begin{align}
\begin{cases} 
\displaystyle \frac{\partial u}{\partial x} = \frac{\partial v}{\partial y}, \\
\displaystyle \frac{\partial u}{\partial y} = -\frac{\partial v}{\partial x},
\end{cases} \label{eq:2.2.10}
\end{align}
就称 $v$ 为 $u$ 的\textbf{共轭调和函数}.
\end{definition}

\begin{proposition}
全纯函数的实部和虚部就构成一对共轭调和函数.
\end{proposition}
\begin{proof}
由\refthe{theorem:全纯函数的实部和虚部都是调和函数}和\refthe{theorem:复变函数可微的充要条件2}立得.
\end{proof}

\begin{theorem}
设 $u$ 是单连通区域 $D$ 上的调和函数，则必存在 $u$ 的共轭调和函数 $v$，使得 $u + \mathrm{i}v$ 是 $D$ 上的全纯函数.
\end{theorem}
\begin{proof}
因为 $u$ 满足 Laplace 方程
\begin{align*}
\frac{\partial^2 u}{\partial x^2} + \frac{\partial^2 u}{\partial y^2} &= 0,
\end{align*}
若令 $P = -\frac{\partial u}{\partial y}, Q = \frac{\partial u}{\partial x}$，则
\begin{align*}
\frac{\partial Q}{\partial x} &= \frac{\partial^2 u}{\partial x^2} = -\frac{\partial^2 u}{\partial y^2} = \frac{\partial P}{\partial y},
\end{align*}
所以
\begin{align*}
P \mathrm{d}x + Q \mathrm{d}y &= -\frac{\partial u}{\partial y} \mathrm{d}x + \frac{\partial u}{\partial x} \mathrm{d}y
\end{align*}
是一个全微分，因而积分
\begin{align*}
\int_{(x_0, y_0)}^{(x, y)} -\frac{\partial u}{\partial y} \mathrm{d}x + \frac{\partial u}{\partial x} \mathrm{d}y
\end{align*}
与路径无关. 令
\begin{align*}
v(x, y) = \int_{(x_0, y_0)}^{(x, y)} -\frac{\partial u}{\partial y} \mathrm{d}x + \frac{\partial u}{\partial x} \mathrm{d}y,
\end{align*}
则
\begin{align*}
v(x,y)&=\int_{(x_0,y_0)}^{(x,y_0)}{-\frac{\partial u}{\partial y}\mathrm{d}x+\frac{\partial u}{\partial x}\mathrm{d}y}+\int_{(x,y_0)}^{(x,y)}{-\frac{\partial u}{\partial y}\mathrm{d}x+\frac{\partial u}{\partial x}\mathrm{d}y}
\\
&=\int_{x_0}^x{-\frac{\partial u}{\partial y}\mathrm{d}x}+0+0+\int_{y_0}^y{\frac{\partial u}{\partial x}\mathrm{d}y}
\\
&=\int_{x_0}^x{-\frac{\partial u}{\partial y}\mathrm{d}x}+\int_{y_0}^y{\frac{\partial u}{\partial x}\mathrm{d}y}.
\end{align*}
那么
\begin{align*}
\begin{cases} 
\displaystyle \frac{\partial v}{\partial x} = -\frac{\partial u}{\partial y}, \\
\displaystyle \frac{\partial v}{\partial y} = \frac{\partial u}{\partial x}.
\end{cases}
\end{align*}
所以，$v$ 就是要求的 $u$ 的共轭调和函数.
\end{proof}

\begin{theorem}[函数在一点解析的充要条件]\label{theorem:函数在一点解析的充要条件}
设$f(z)=u(x,y)+iv(x,y)$在点$z_0$的某个邻域内有定义, 则$f(z)$在$z_0$点解析的充要条件是以下条件之一:
\begin{enumerate}[(1)]
\item $f(z)$在$z_0$的某个邻域内处处可导.
\item $u(x,y)$和$v(x,y)$在$z_0$的某个邻域内具有连续的一阶偏导数, 且在该邻域内满足Cauchy-Riemann方程:
\begin{align*}
\frac{\partial u}{\partial x} = \frac{\partial v}{\partial y}, \quad \frac{\partial v}{\partial x} = -\frac{\partial u}{\partial y}.
\end{align*}
\item $f(z)$在$z_0$的某个邻域内可以展开成$(z-z_0)$的幂级数, 即
\begin{align*}
f(z) = \sum\limits_{n=0}^{\infty} c_n (z-z_0)^n.
\end{align*}
\item $v(x,y)$在$z_0$的某个邻域内是$u(x,y)$的共轭调和函数.
\end{enumerate}
\end{theorem}

\begin{theorem}[函数在区域内解析的充要条件]\label{theorem:函数在区域内解析的充要条件}
设$D\subset \mathbb{C}$,则
\begin{enumerate}[(1)]
\item 函数 \(f(z) = f(x,y) = u(x,y) + \mathrm{i}v(x,y)\) 在区域 \(D\) 内解析的充要条件是\(u, v\) 在$D$内是可微的，且满足 Cauchy-Riemann 方程：
\begin{align*}
\frac{\partial u}{\partial x} = \frac{\partial v}{\partial y}, \quad \frac{\partial v}{\partial x} = -\frac{\partial u}{\partial y}.
\end{align*}

\item 函数 \(f(z) = f(x,y) = u(x,y) + \mathrm{i}v(x,y)\) 在 \(D\) 内解析的充要条件为：
\begin{enumerate}[(i)]
\item \(u_x, u_y, v_x, v_y\) 在 \(D\) 内连续；
\item \(u(x,y), v(x,y)\) 在 \(D\) 内满足Cauchy-Riemann方程.
\end{enumerate}

\item 函数 \(f(z)\) 在区域 \(D\) 内解析的充要条件是：
\begin{enumerate}[(i)]
\item \(f(z)\) 在 \(D\) 内连续.
\item 对任一可求长简单闭曲线\(C\)，只要 \(C\) 及其内部全含于 \(D\) 内，就有 \(\int_{C} f(z) \mathrm{d}z = 0\).
\end{enumerate}

\item 函数 \(f(z) = f(x,y) = u(x,y) + \mathrm{i}v(x,y)\) 在 \(D\) 内解析的充要条件为在区域 \(D\) 内 \(v(x,y)\) 是 \(u(x,y)\) 的共轭调和函数.

\item 函数 \(f(z)\) 在 \(D\) 内解析的充要条件为\(f(z)\) 在 \(D\) 内可局部(逐点)展开为幂级数,即$f$在$D$内任一点 \(a\) 的邻域内可展开成 \(z - a\) 的幂级数,, 即
\begin{align*}
f(z) = \sum\limits_{n=0}^{\infty} c_n (z-z_0)^n.
\end{align*}
\end{enumerate}
\end{theorem}

\begin{example}
设$f(z)=|z|^2$, 则以下常见错误命题均不成立:
\begin{enumerate}[(1)]
\item 若$f(z)$在$z_0$不解析则不可导. 反例: $f(z)=|z|^2$在$z=0$处可导但不解析.
\item 若$f(z)$在$z_0$可导则解析. 反例: $f(z)=|z|^2$在$z=0$处可导但不解析.
\item 若$u,v$在$z_0$满足C-R方程则$f(z)$在$z_0$解析. 反例: $f(z)=|z|^2$在$z=0$处满足C-R方程, 但在$z=0$的任意邻域内不满足C-R方程, 因此不解析.
\end{enumerate}
\end{example}






\end{document}